% Curated XeLaTeX preamble for academic/research markdown reports.
% Tuned for a *modern web* look: sans-serif body, GitHub-style code blocks,
% colored hyperlinks, generous spacing. Loaded by pandoc via --include-in-header.

% ---- Math ---------------------------------------------------------------
\usepackage{amsmath}
\usepackage{amssymb}
\usepackage{mathtools}
\usepackage{unicode-math}
% STIX Two Math pairs cleanly with sans-serif body fonts and gives wide
% Unicode coverage. Falls back to Latin Modern Math if STIX absent.
\IfFontExistsTF{STIX Two Math}{%
  \setmathfont{STIX Two Math}%
}{%
  \setmathfont{Latin Modern Math}%
}

% ---- CJK (Chinese / Japanese / Korean) ---------------------------------
% Load xeCJK only when a usable CJK font is present. Falls through silently
% on systems without CJK fonts so pure-Latin builds are untouched.
% Order: macOS system fonts → Noto CJK (Linux/CTAN) → give up.
\IfFontExistsTF{PingFang SC}{%
  \usepackage{xeCJK}%
  \setCJKmainfont[BoldFont={PingFang SC Semibold}, ItalicFont={Kaiti SC}]{PingFang SC}%
  \IfFontExistsTF{Heiti SC}{\setCJKsansfont{Heiti SC}}{}%
  \IfFontExistsTF{STSong}{\setCJKmonofont{STSong}}{}%
}{%
  \IfFontExistsTF{Noto Serif CJK SC}{%
    \usepackage{xeCJK}%
    \setCJKmainfont{Noto Serif CJK SC}%
    \IfFontExistsTF{Noto Sans CJK SC}{\setCJKsansfont{Noto Sans CJK SC}}{}%
    \IfFontExistsTF{Noto Sans Mono CJK SC}{\setCJKmonofont{Noto Sans Mono CJK SC}}{}%
  }{%
    \IfFontExistsTF{Songti SC}{%
      \usepackage{xeCJK}%
      \setCJKmainfont{Songti SC}%
    }{}%
  }%
}

% ---- Unicode symbol fallbacks ------------------------------------------
% Some text fonts lack glyphs for Greek letters and math symbols that appear
% *outside* math mode in research prose (e.g. "C(α) 检验", "x → y", "p × q").
% Without these mappings such characters silently vanish in the PDF.
% Routing them through math mode lets the math font render them.
\usepackage{newunicodechar}
% Greek lowercase
\newunicodechar{α}{\ensuremath{\alpha}}
\newunicodechar{β}{\ensuremath{\beta}}
\newunicodechar{γ}{\ensuremath{\gamma}}
\newunicodechar{δ}{\ensuremath{\delta}}
\newunicodechar{ε}{\ensuremath{\varepsilon}}
\newunicodechar{ζ}{\ensuremath{\zeta}}
\newunicodechar{η}{\ensuremath{\eta}}
\newunicodechar{θ}{\ensuremath{\theta}}
\newunicodechar{ι}{\ensuremath{\iota}}
\newunicodechar{κ}{\ensuremath{\kappa}}
\newunicodechar{λ}{\ensuremath{\lambda}}
\newunicodechar{μ}{\ensuremath{\mu}}
\newunicodechar{ν}{\ensuremath{\nu}}
\newunicodechar{ξ}{\ensuremath{\xi}}
\newunicodechar{π}{\ensuremath{\pi}}
\newunicodechar{ρ}{\ensuremath{\rho}}
\newunicodechar{σ}{\ensuremath{\sigma}}
\newunicodechar{τ}{\ensuremath{\tau}}
\newunicodechar{υ}{\ensuremath{\upsilon}}
\newunicodechar{φ}{\ensuremath{\varphi}}
\newunicodechar{χ}{\ensuremath{\chi}}
\newunicodechar{ψ}{\ensuremath{\psi}}
\newunicodechar{ω}{\ensuremath{\omega}}
% Greek uppercase (only the ones that differ from Latin)
\newunicodechar{Γ}{\ensuremath{\Gamma}}
\newunicodechar{Δ}{\ensuremath{\Delta}}
\newunicodechar{Θ}{\ensuremath{\Theta}}
\newunicodechar{Λ}{\ensuremath{\Lambda}}
\newunicodechar{Ξ}{\ensuremath{\Xi}}
\newunicodechar{Π}{\ensuremath{\Pi}}
\newunicodechar{Σ}{\ensuremath{\Sigma}}
\newunicodechar{Φ}{\ensuremath{\Phi}}
\newunicodechar{Ψ}{\ensuremath{\Psi}}
\newunicodechar{Ω}{\ensuremath{\Omega}}
% Arrows
\newunicodechar{→}{\ensuremath{\rightarrow}}
\newunicodechar{←}{\ensuremath{\leftarrow}}
\newunicodechar{↔}{\ensuremath{\leftrightarrow}}
\newunicodechar{⇒}{\ensuremath{\Rightarrow}}
\newunicodechar{⇐}{\ensuremath{\Leftarrow}}
\newunicodechar{⇔}{\ensuremath{\Leftrightarrow}}
\newunicodechar{↦}{\ensuremath{\mapsto}}
% Operators / relations
\newunicodechar{×}{\ensuremath{\times}}
\newunicodechar{÷}{\ensuremath{\div}}
\newunicodechar{±}{\ensuremath{\pm}}
\newunicodechar{∓}{\ensuremath{\mp}}
\newunicodechar{·}{\ensuremath{\cdot}}
\newunicodechar{≤}{\ensuremath{\leq}}
\newunicodechar{≥}{\ensuremath{\geq}}
\newunicodechar{≠}{\ensuremath{\neq}}
\newunicodechar{≈}{\ensuremath{\approx}}
\newunicodechar{≡}{\ensuremath{\equiv}}
\newunicodechar{∼}{\ensuremath{\sim}}
\newunicodechar{∝}{\ensuremath{\propto}}
\newunicodechar{∞}{\ensuremath{\infty}}
% Set theory / logic
\newunicodechar{∈}{\ensuremath{\in}}
\newunicodechar{∉}{\ensuremath{\notin}}
\newunicodechar{∋}{\ensuremath{\ni}}
\newunicodechar{⊂}{\ensuremath{\subset}}
\newunicodechar{⊆}{\ensuremath{\subseteq}}
\newunicodechar{⊃}{\ensuremath{\supset}}
\newunicodechar{⊇}{\ensuremath{\supseteq}}
\newunicodechar{∪}{\ensuremath{\cup}}
\newunicodechar{∩}{\ensuremath{\cap}}
\newunicodechar{∅}{\ensuremath{\emptyset}}
\newunicodechar{∀}{\ensuremath{\forall}}
\newunicodechar{∃}{\ensuremath{\exists}}
\newunicodechar{¬}{\ensuremath{\neg}}
\newunicodechar{∧}{\ensuremath{\wedge}}
\newunicodechar{∨}{\ensuremath{\vee}}
% Calculus / non-large operators
% NOTE: do NOT add \newunicodechar entries for variable-size / "large"
% operators (∑ ∏ ∫ ∮ ⋃ ⋂ ∐ ⨁ ⨂ etc.). Under unicode-math, macros like
% \sum emit the literal Unicode codepoint back into the math token stream;
% making that codepoint an active char that re-expands to \sum causes
% infinite recursion and xelatex hangs. If a user truly needs a bare ∑
% in prose, write it as $\sum$ in math mode instead of as a literal char.
\newunicodechar{∂}{\ensuremath{\partial}}
\newunicodechar{∇}{\ensuremath{\nabla}}
\newunicodechar{√}{\ensuremath{\surd}}
% Common typography that some text fonts still miss
\newunicodechar{…}{\ensuremath{\ldots}}
\newunicodechar{ℝ}{\ensuremath{\mathbb{R}}}
\newunicodechar{ℕ}{\ensuremath{\mathbb{N}}}
\newunicodechar{ℤ}{\ensuremath{\mathbb{Z}}}
\newunicodechar{ℚ}{\ensuremath{\mathbb{Q}}}
\newunicodechar{ℂ}{\ensuremath{\mathbb{C}}}
\newunicodechar{𝔼}{\ensuremath{\mathbb{E}}}
\newunicodechar{ℙ}{\ensuremath{\mathbb{P}}}

% ---- Color palette (web-inspired) --------------------------------------
\usepackage{xcolor}
\definecolor{headingink}{HTML}{111827}     % near-black for headings (slate-900)
\definecolor{bodyink}{HTML}{1F2937}         % slate-800 for body text
\definecolor{linkink}{HTML}{0969DA}         % GitHub blue for links
\definecolor{mutedink}{HTML}{57606A}        % slate-500 for muted text
\definecolor{codebg}{HTML}{F6F8FA}          % GitHub code background
\definecolor{coderule}{HTML}{D0D7DE}        % subtle border on code blocks
\color{bodyink}                             % default body color

% ---- Tables / floats ----------------------------------------------------
\usepackage{booktabs}
\usepackage{longtable}
\usepackage{array}
\renewcommand{\arraystretch}{1.25}          % more breathing room per row

% ---- Lists --------------------------------------------------------------
\usepackage{enumitem}
\setlist{itemsep=4pt, topsep=6pt, parsep=0pt, leftmargin=1.4em}

% ---- Spacing ------------------------------------------------------------
\usepackage{parskip}                        % blank-line paragraphs, no first-line indent
\usepackage{setspace}
\setstretch{1.4}                            % web-like line spacing

% ---- Section headings ---------------------------------------------------
\usepackage{titlesec}
\titleformat*{\section}{\Huge\bfseries\color{headingink}}
\titleformat*{\subsection}{\LARGE\bfseries\color{headingink}}
\titleformat*{\subsubsection}{\Large\bfseries\color{headingink}}
\titleformat*{\paragraph}{\large\bfseries\color{headingink}}
\titlespacing*{\section}{0pt}{24pt}{10pt}
\titlespacing*{\subsection}{0pt}{18pt}{8pt}
\titlespacing*{\subsubsection}{0pt}{14pt}{6pt}

% ---- Code ---------------------------------------------------------------
\usepackage{fvextra}                        % adds breaklines to fancyvrb
\usepackage{framed}
% Pandoc's default Shaded env wraps Highlighting blocks. Retint with the
% GitHub palette and bump internal padding so code looks airy, not cramped.
\definecolor{shadecolor}{HTML}{F6F8FA}
\setlength{\FrameSep}{10pt}
\setlength{\OuterFrameSep}{6pt}
% Enable line wrapping for both plain Verbatim (unhighlighted code) and the
% Highlighting environment pandoc emits for fenced code blocks. Without
% breakanywhere, long identifiers/URLs run off the right margin.
\fvset{breaklines=true,breakanywhere=true,fontsize=\small,breaksymbolleft={},breaksymbolright={}}
\DefineVerbatimEnvironment{Highlighting}{Verbatim}{commandchars=\\\{\},breaklines,breakanywhere,fontsize=\small,breaksymbolleft={},breaksymbolright={}}

% ---- Hyperlinks (must be near the end) ----------------------------------
\usepackage{hyperref}
\hypersetup{
  colorlinks=true,
  linkcolor=linkink,
  citecolor=linkink,
  urlcolor=linkink,
  filecolor=linkink,
  pdfborder={0 0 0},
  bookmarksnumbered=true,
  bookmarksopen=true,
  pdfstartview={FitH}
}

% ---- Footnotes ----------------------------------------------------------
\usepackage[bottom,hang,flushmargin]{footmisc}

% ---- Quotes -------------------------------------------------------------
\usepackage{csquotes}

% ---- Title block tightening --------------------------------------------
\makeatletter
\renewcommand{\maketitle}{%
  \begingroup
    \centering
    {\Huge\bfseries\color{headingink} \@title \par}\vspace{0.8em}%
    \ifx\@author\@empty\else{\large\color{mutedink} \@author \par}\vspace{0.4em}\fi
    \ifx\@date\@empty\else{\small\color{mutedink} \@date \par}\fi
    \vspace{1.6em}%
  \endgroup
}
\makeatother
